\chapter{Σχεδιασμός και Υλοποίηση Αισθητήρα}
\label{ch:sensor}

\section{Εισαγωγή}

Στο συγκεκριμένο κεφάλαιο θα εξεταστεί αναλυτικά η αρχιτεκτονική και οι
αρχές λειτουργίας του συστήματος. Παρουσιάζονται τα δομικά στοιχεία με
τα χαρακτηριστικά τους, η θεωρία λειτουργίας του αισθητήρα, το
σχηματικό, οι απαιτήσεις και οι περιορισμοί του συστήματος καθώς και η
τελική πλακέτα.

\section{Δομικά Στοιχεία του Συστήματος}

\subsection{Γεννήτρια DDS}

Ο υπολογισμός της σύνθετης αντίστασης του δείγματος απαιτεί τη διέγερση
του από ένα ηλεκτρικό σήμα μέσω των ηλεκτροδίων και μίας αντίστασης
(shunt resistor) συνδεδεμένη σε σειρά. Για τη δημιουργία του σήματος
επιλέχθηκε ο DDS (Direct Digital Synthesis) AD9833. Πρόκειται για ένα
ολοκληρωμένο κύκλωμα που επιτρέπει την παραγωγή ημιτοειδών, τετραγωνικών
και τριγωνικών κυματομορφών. Επιλέχθηκε λόγω του χαμηλού κόστους και
κατανάλωσης ισχύος, αλλά και τον εύκολο έλεγχο μέσω SPI για τον
προγραμματισμό του επιθυμητού σήματος και συχνότητας μέσω του
μικροελεγκτή. Υποστηρίζει συχνότητες εξόδου από 0-12.5 MHz και παρέχει
πλάτος εξόδου περίπου 0,6Vpp. \cite{ref32}

\begin{figure}[H]
  \centering
  \includegraphics[width=0.34\textwidth]{ad9833-dds-module.png}
  \caption{Μονάδα γεννήτριας DDS AD9833.}
  \label{fig:ad9833-module}
\end{figure}

\subsection{Ανιχνευτής Κέρδους-Φάσης}

Για τον υπολογισμό της πτώσης τάσης και της διαφοράς φάσης μεταξύ τάσης
και ρεύματος του δείγματος χρησιμοποιήθηκε η μονάδα \textbf{AD8302}.
Πρόκειται για ένα ολοκηρωμένο κύκλωμα ανίχνευσης κέρδους και φάσης
μεταξύ δύο AC σημάτων εισόδου, το οποίο παράγει στην έξοδο δύο DC τιμές
τάσεων, μία για το λόγο των πλατών και μία για τη διαφορά φάσης
αντίστοιχα. Η ευαισθησία της εξόδου είναι 30mV/db για το κέρδος και
10mV/° για τη φάση, ενώ μπορεί να δεχθεί σήματα συχνότητας έως και
2.7GHz. Παρακάτω φαίνεται το εξιδανικευμένο διάγραμμα εισόδων-εξόδου για
την εκάστοτε έξοδο. \cite{ref33}

\begin{figure}[H]
  \centering
  \includegraphics[width=0.76\textwidth]{ad8302-phase-transfer.png}
  \caption{Ιδανική χαρακτηριστική εξόδου φάσης του AD8302.}
  \label{fig:ad8302-phase-transfer}
\end{figure}

\begin{figure}[H]
  \centering
  \includegraphics[width=0.76\textwidth]{ad8302-magnitude-transfer.png}
  \caption{Ιδανική χαρακτηριστική εξόδου λόγου πλατών του AD8302.}
  \label{fig:ad8302-magnitude-transfer}
\end{figure}

Η συνάρτηση μεταφοράς της εξόδου για το λόγο των πλατών δίνεται από τη
σχέση

\[V_{MAG} = 0,9 + 0,03 \cdot \left( 20\log\frac{\left| V_{a} \right|}{\left| V_{b} \right|} \right)\]

ενώ για τη φάση ισχύει

\[V_{PHASE} = 0,9 - 0,01 \cdot \left( \Delta\varphi - 90^{\circ} \right)\]

\subsection{Μονάδα Analog-to-Digital Converter}

Για τη δειγματοληψία των τιμών εξόδου του ανιχνευτή κέρδους-φάσης
χρησιμοποιήθηκε o εξωτερικός μετατροπέας αναλογικών σε ψηφιακό σήμα
(ADC) \textbf{ADS1115} των 16bits για καλύτερη ακρίβεια. Διαθέτει
τέσσερα κανάλια εισόδου κι επικοινωνεί με το πρωτόκολλο
I\textsuperscript{2}C. \cite{ref34}

\begin{figure}[H]
  \centering
  \includegraphics[width=0.42\textwidth]{ads1115-adc-module.jpg}
  \caption{Μονάδα μετατροπέα αναλογικού σε ψηφιακό σήμα ADS1115 των 16bits.}
  \label{fig:ads1115-module}
\end{figure}

\subsection{Ολοκληρωμένα Κυκλώματα}

Για τη διαμόρφωση των αναλογικών σημάτων (signal conditioning)
επιλέχθηκε ο \textbf{AD8066} λόγω της χαμηλής του κατανάλωσης. Για την
απομόνωση των σημάτων του κυκλώματος χρησιμοποιήθηκε ο \textbf{LTC6268}
σε συνδεσμολογία ακόλουθου τάσης (voltage follower). Η επιλογή του
βασίστηκε στην εξαιρετικά χαμηλή χωρητικότητα εισόδου της τάξης των fF
και στο πολύ χαμηλό ρεύμα πόλωσης (input bias current) στην κλίμακα των
pA, προκειμένου να ελαχιστοποιηθεί η φόρτωση του δείγματος και να
αποφευχθεί η αλλοίωση των πραγματικών ηλεκτρικών χαρακτηριστικών. \cite{ref35,ref36}

Στη συνέχεια, τα παραπάνω σήματα οδηγούνται μέσω του αναλογικού διακόπτη
\textbf{ADG733} στους διαφορικούς ενισχυτές \textbf{AD8130.} Πρόκειται
για έναν high-speed διαφορικό ενισχυτή που δεν απαιτεί εξωτερικά
εξαρτήματα και προσφέρει υψηλή απόρριψη κοινού θορύβου, με περίπου -80db
στο 1MHz. \cite{ref37,ref38}

\subsection{Μικροελεγτκής}

Για τον έλεγχο και το συντονισμό της λειτουργίας του συστήματος
χρησιμοποιήθηκε ο μικροελεγκτής ESP32-NodeMCU-32s. Ο μικροελεγκτής
αναλαμβάνει το ρόλο του προγραμματισμού της γεννήτριας AD9833 μέσω του
πρωτοκόλλου SPI, τη λήψη των δεδομένων από τον μετατροπέα ADS1115 με το
πρωτόκολλο \(I^{2}C\), τον έλεγχο του αναλογικού διακόπτη και τέλος την
εκτέλεση των απαραίτητων πράξεων για τον υπολογισμό της σύνθετης
αντίστασης και την εκτίμηση της υγρασίας του δείγματος. \cite{ref39}

\section{Απαιτήσεις του Συστήματος}

Μία από τις απαιτήσεις του συστήματος είναι η ανίχνευση μικρών μεταβολών
της χωρητικότητας. Σύμφωνα με τις μετρήσεις που πραγματοποιήθηκαν με τον
Αναλυτή Εμπέδησης, η διακύμανση της χωρητικότητας στα δείγματα που
εξετάστηκαν ήταν της τάξης των δεκάδων pF. Τόσο χαμηλές χωρητικότητες
μπορούν εύκολα να επισκιαστούν από παρασιτικές χωρητικότητες,
παρεμβολές, ακατάλληλο σχεδιασμό της πλακέτας και γενικότερα συστημικά
σφάλματα. Έτσι, δόθηκε ιδιαίτερη προσοχή στο σχεδιασμό προκειμένου να
ελαχιστοποιηθούν τα σφάλματα στο μέγιστο δυνατό.

Ακόμη, ένα μειονέκτημα του ανιχνευτή κέρδους-φάσης είναι η μη
γραμμικότητα των συναρτήσεων εξόδου για διαφορά φάσης από
0\textsuperscript{ο} έως +-30\textsuperscript{ο} (\figref{fig:ad8302-phase-error}). Από τον
αναλυτή εμπέδησης φάνηκε πως η φάση των δειγμάτων κυμαινόταν από τιμές
κοντά στις 0\textsuperscript{ο} έως και τις -10\textsuperscript{ο}. Για
αυτό το λόγο, δημιουργήθηκε ένα τρίτο σήμα αναφοράς Vereference με
αρκετά μεγαλύτερη διαφορά φάσης συγκριτικά με τα δύο σήματα που
λαμβάνονται από το σύστημα, προκειμένου να μειωθεί το σφάλμα εξόδου.
Εφόσον ληφθεί η μέτρηση του κάθε σήματος αναφορικά με το Vreference,
γίνεται μαθηματικά η απαλοιφή του Vreference μέσω του μικροελεγκτή και
εξάγεται το τελικό V\textsubscript{MAG} και V\textsubscript{PHASE} των
δύο σημάτων ενδιαφέροντος.

\begin{figure}[H]
  \centering
  \includegraphics[width=0.70\textwidth]{ad8302-phase-error.png}
  \caption{Χαρακτηριστική εξόδου φάσης και σφάλματος του AD8302 συναρτήσει της διαφοράς φάσης. \cite{ref33}}
  \label{fig:ad8302-phase-error}
\end{figure}

\section{Προσέγγιση και Μεθοδολογία}

Στο \figref{fig:sensor-architecture} φαίνεται η συνολική αρχιτεκτονική του συστήματος. Αρχικά,
ο DDS παράγει ένα ημιτονοειδές σήμα συχνότητας 1MHz και πλάτους 0,6mVpp,
το οποίο περιλαμβάνει μία DC συνιστώσα (DC offset). Ωστόσο, επειδή μία
DC συνιστώσα θα προκαλούσε οξείδωση στα ηλεκτρόδια λόγω της
ηλεκτρόλυσης, το σήμα θα πρέπει να είναι καθαρά εναλλασσόμενο (AC). Για
αυτό το λόγο, στην έξοδο του DDS έχει τοποθετηθεί ένα υψηπερατό φίλτρο
όπως φαίνεται στο \figref{fig:sensor-architecture} με συχνότητα αποκοπής

\[f = \frac{1}{2\pi * R_{21} * C_{34}} = \frac{1}{2\pi * 1k * 0.1u} = 1.6kHz\]

\begin{figure}[H]
  \centering
  \includegraphics[width=0.96\textwidth]{sensor-system-architecture.png}
  \caption{Συνολική αρχιτεκτονική του συστήματος.}
  \label{fig:sensor-architecture}
\end{figure}

\begin{figure}[H]
  \centering
  \includegraphics[width=0.72\textwidth]{dds-high-pass-filter.png}
  \caption{Η μονάδα παραγωγής σημάτων και το υψηπερατό φίλτρο}
  \label{fig:dds-high-pass-filter}
\end{figure}

Στη συνέχεια, για τη βελτίωση του SNR το σήμα ενισχύεται μέσω ενός
τελεστικού σε συνδεσμολογία αναστρέφοντα ενισχυτή με κέρδος

\[A_{v} = - \ \frac{R_{14}}{R_{13}} = \  - \ \frac{10k}{6.8k} = - 1.5\]

Το σήμα αυτό χρησιμοποιείται για τη διέγερση του δείγματος αλλά και την
παραγωγή του σήματος αναφοράς, οπότε και στέλνεται μέσω του All Pass
φίλτρου για τη δημιουργία καθυστέρησης φάσης περίπου \(35^{\circ}\) σε σχέση
με το αρχικό σήμα (\figref{fig:excitation-reference-conditioning}).

\begin{figure}[H]
  \centering
  \includegraphics[width=0.88\textwidth]{excitation-reference-conditioning.png}
  \caption{Η διαμόρφωση των σημάτων διέγερσης και αναφοράς. Ο άνω τελεστικός λειτουργεί ως αναστρέφων ενισχυτής ενώ ο κάτω τελεστικός σε συνδεσμολογία All Pass Filter εισάγει τη διαφορά φάσης.}
  \label{fig:excitation-reference-conditioning}
\end{figure}

Κατά την εφαρμογή του σήματος διέγερσης στο δείγμα και την αντίσταση
αναφοράς (R10), τα σήματα στα άκρα των δύο στοιχείων λαμβάνονται μέσω
κυκλωμάτων απομόνωσης (buffers) τα οποία διαθέτουν πολύ υψηλή αντίσταση
αλλά χαμηλή χωρητικότητα εισόδου (\figref{fig:signal-buffering-filtering}). Στη συνέχεια, οδηγούνται
στους διαφορικούς ενισχυτές για τον υπολογισμό της πτώσης τάσεων στα δύο
στοιχεία (\figref{fig:differential-amplifiers-switching}).

\begin{figure}[H]
  \centering
  \includegraphics[width=0.82\textwidth]{signal-buffering-filtering.png}
  \caption{Οι τελεστικοί στα άκρα της αντίστασης και του δείγματος σε συνδεσμολογία ακολούθου τάσης και τα χαμηλοπερατά φίλτρα στις εξόδους τους.}
  \label{fig:signal-buffering-filtering}
\end{figure}

Τα δύο σήματα εξόδου των διαφορικών ενισχυτών, οδηγούνται σε έναν
αναλογικό διακόπτη, ο οποίος ελέγχεται από μία έξοδο (GPIO) του ESP32.
Τελικά, στην είσοδο του ανιχνευτή κέρδους φάσης καταλήγει το σήμα
αναφοράς και, εναλλάξ, τα δύο σήματα που επιλέγονται από τον
διακόπτη προκειμένου να πραγματοποιηθεί η μέτρηση για το κάθε στοιχείο
ξεχωριστά (\figref{fig:differential-amplifiers-switching}). Η έξοδος του ανιχνευτή δειγματοληπτείται από τον
εξωτερικό μετατροπέα αναλογικού σε ψηφιακό ώστε να γίνουν οι απαραίτητοι
υπολογισμοί μέσω του μικροελεγκτή για την εξαγωγή των ηλεκτρικών
παραμέτρων του δείγματος (\figref{fig:phase-gain-detector-adc})

\begin{figure}[H]
  \centering
  \includegraphics[width=0.7\textwidth]{differential-amplifiers-switching.png}
  \caption{Τα σήματα οδηγούνται στους  διαφορικικούς ενισχυτές όπου μέσω δύο χαμηλοπερατών φίλτρων παιρνούν στον αναλογικό διακόπτη.}
  \label{fig:differential-amplifiers-switching}
\end{figure}

\begin{figure}[H]
  \centering
  \includegraphics[width=0.7\textwidth]{phase-gain-detector-adc.png}
  \caption{Τα σήματα οδηγούνται στον ανιχνευτή κέρδους-τάσης και στη συνέχεια οι έξοδοι του διαβάζονται από τον ADC.}
  \label{fig:phase-gain-detector-adc}
\end{figure}

Έστω ότι ορίζουμε ως \(V_{A}\) τη διαφορά τάσης στα άκρα των
ηλεκτροδίων, \(V_{B}\) τη διαφορά τάσης στα άκρα της αντίστασης και
\(V_{ref}\) το διαμορφωμένο σήμα με καθυστέρηση φάσης \(35^{\circ}\). Ως \(\left|\Delta\varphi_{1}\right|\) ορίζεται το μέτρο της διαφοράς φάσης μεταξύ  \(V_{A}\) και του σήματος αναφοράς \(V_{ref}\) . Αντίστοιχα, ως \(\left|\Delta\varphi_{2}\right|\) ορίζεται το μέτρο της διαφοράς φάσης μεταξύ  \(V_{B}\) και του \(V_{ref}\) . Μετά
από δύο εναλλαγές του διακόπτη, μετρώνται τα παρακάτω ζεύγη μεγεθών

\[
V_{\mathrm{MAG}1}=0{,}9+0{,}03\left(20\log_{10}
\frac{\left|V_{A}\right|}{\left|V_{ref}\right|}\right)
\tag{10}
\]

\[
V_{\mathrm{PHASE}1}
=
0{,}9
-
0{,}01
\left(
\left|\Delta\varphi_{1}\right|
-
90^{\circ}
\right)
\tag{11}
\]

και

\[
V_{\mathrm{MAG}2}=0{,}9+0{,}03\left(20\log_{10}
\frac{\left|V_{B}\right|}{\left|V_{ref}\right|}\right)
\tag{12}
\]

\[
V_{\mathrm{PHASE}2}
=
0,9
-
0,01
\left(
\left|\Delta\varphi_{2}\right|
-
90^{\circ}
\right)
\tag{13}
\]

Για τη σχέση (10), ισχύει

\[
\log_{10}\left(\frac{\left|V_{A}\right|}{\left|V_{ref}\right|}\right)
=\frac{V_{\mathrm{MAG}1}-0{,}9}{0{,}6}
\]

και επομένως

\[
\left|V_{A}\right|=
10^{\frac{V_{\mathrm{MAG}1}-0{,}9}{0{,}6}}\left|V_{ref}\right|
\tag{14}
\]

Ενώ αντίστοιχα για το \(V_{B}\) ισχύει

\[
\left|V_{B}\right|=
10^{\frac{V_{\mathrm{MAG}2}-0{,}9}{0{,}6}}\left|V_{ref}\right|
\tag{15}
\]

Για να βρεθεί ο λόγος των πλατών, διαιρούνται κατά μέλη οι σχέσεις (14)
και (15) και προκύπτει τελικά

\[
\frac{\left|V_{A}\right|}{\left|V_{B}\right|}
=10^{\frac{V_{\mathrm{MAG}1}-V_{\mathrm{MAG}2}}{0{,}6}}
\tag{16}
\]

Από τις σχέσεις (11) και (13) υπολογίζονται αντίστοιχα οι διαφορές φάσης

\[
\left|\Delta\varphi_{1}\right|
=\frac{0{,}9-V_{\mathrm{PHASE}1}}{0{,}01}+90^{\circ}
\tag{17}
\]

και

\[
\left|\Delta\varphi_{2}\right|
=\frac{0{,}9-V_{\mathrm{PHASE}2}}{0{,}01}+90^{\circ}.
\tag{18}
\]

Η διαφορά φάσης των δύο σημάτων ενδιαφέροντος προκύπτει με απαλοιφή της
κοινής φάσης του σήματος αναφοράς. Θεωρώντας ότι το σήμα αναφοράς προηγείται κατά $35^\circ$ περίπου μοίρες από τα άλλα δύο, έχουμε $\varphi_1-\varphi_3>0$ και $\varphi_2-\varphi_3>0$ όπου $\varphi_3$ το σήμα αναφοράς, επομένως ισχύει

\[
\Delta\varphi
= |\Delta\varphi_1|-|\Delta\varphi_2|
= |\varphi_1-\varphi_3|-|\varphi_2-\varphi_3|
= \Delta\varphi_1-\Delta\varphi_2
= \frac{V_{\mathrm{PHASE}2}-V_{\mathrm{PHASE}1}}{0{,}01}.
\tag{19}
\]

Θεωρώντας \(R\) την αντίσταση αναφοράς \(R_{10}\), το ίδιο ρεύμα διαρρέει
την αντίσταση και το δείγμα. Επομένως, η σύνθετη αντίσταση του δείγματος
υπολογίζεται ως

\[
\begin{aligned}
\vec{Z}
&=\frac{\vec{V}_{A}}{\vec{I}}
 =\frac{\vec{V}_{A}}{\vec{V}_{B}/R}
 =R\frac{\vec{V}_{A}}{\vec{V}_{B}} \\
&=R\frac{\left|V_{A}\right|}{\left|V_{B}\right|}
  \angle\left(\Delta\varphi_{1}-\Delta\varphi_{2}\right) \\
&=R\,10^{\frac{V_{\mathrm{MAG}1}-V_{\mathrm{MAG}2}}{0{,}6}}
  \angle\left(\frac{V_{\mathrm{PHASE}2}-V_{\mathrm{PHASE}1}}{0{,}01}\right).
\end{aligned}
\tag{20}
\]

Για την εξαγωγή των παραμέτρων του παράλληλου ισοδύναμου
\(R_{p}\)-\(C_{p}\), χρησιμοποιείται η αγωγιμότητα

\[
\begin{aligned}
\vec{Y}
&=\frac{1}{\vec{Z}}
 =\frac{1}{\left|Z\right|e^{j\varphi}} \\
&=\frac{1}{\left|Z\right|}
  \left(\cos\varphi-j\sin\varphi\right).
\end{aligned}
\tag{21}
\]

Από το Κεφάλαιο 3 θεωρώντας ως ηλεκτρικό ισοδύναμο το παράλληλο Rp-Cp, ισχύει επίσης

\[
\vec{Y}=\frac{1}{R_{p}}+j\omega C_{p}=G+jB,
\tag{22}
\]

όπου \(\omega=2\pi f\). Εξισώνοντας το πραγματικό και το φανταστικό μέρος
των σχέσεων (21) και (22), προκύπτουν

\[
G=\frac{1}{R_{p}}=\frac{\cos\varphi}{\left|Z\right|},
\qquad
B=\omega C_{p}=-\frac{\sin\varphi}{\left|Z\right|}.
\tag{23}
\]

Τελικά, οι παράμετροι του ισοδύναμου κυκλώματος δίνονται από τις σχέσεις


\[
R_{p}=\frac{\left|Z\right|}{\cos\varphi},
\qquad
C_{p}=-\frac{\sin\varphi}{\omega\left|Z\right|}.
\tag{24}
\]




\section{Τυπωμένο Κύκλωμα}

Η πλακέτα σχεδιάστηκε με τη χρήση του λογισμικού Kicad. Είναι διπλής όψεος κι έχει διαστάσεις περίπου 10cm x 17cm (Σχήμα 4.12). Τα σήματα διατηρήθηκαν όσο ήταν εφικτό στο άνω επίπεδο, προκειμένου το κάτω επίπεδο που ορίστηκε ως η γείωση να παραμείνει όσο το δυνατό συνεχές. Τα αναλογικά σήματα συγκεντρώθηκαν στην αριστερή μερία της πλακέτας, ενώ τα ψηφιακά τοποθετήθηκαν στα δεξιά για να αποφευχθούν ανεπιθύμητες συζεύξεις (crosstalk). Τα σήματα που οδηγούνται στο διαφορικό ενισχυτή σχεδιάστηκαν ως  διαφορικό ζεύγος (differential pair), με στόχο τη μείωση του κοινού θορύβου. Η μονάδα AD8302 διαθέτει εισόδους SMA, επομένως για τη μεταφορά των σημάτων από την πλακέτα στη μονάδα ανίχνευσης κέρδους-φάσης χρησιμοποιήθηκαν U.FL connectors και καλώδια U.FL-SMA.

\begin{figure}[H]
  \centering
  \includegraphics[width=0.9\textwidth]{SoilBoard.png}
  \caption{Το προσχέδιο της πλακέτας από το λογισμικό σχεδίασης. }
  \label{fig:signal-buffering-filtering}
\end{figure}

Τα ολοκληρωμένα κυκλώματα είναι τύπου SMD, ενώ τα υπόλοιπα στοιχεία είναι through hole.
Οι γραμμές τροφοδοσίας των ολοκληρωμένων κυκλωμάτων είναι στα +2.5V και -2.5V ενώ ο ESP32 τροφοδοτείται με 3.3V. Στις τροφοδοσίες των ολοκληρωμένων τοποθετήθηκαν κεραμικοί πυκνωτές αποσύζευξης (decoupling capacitors) με τιμές 0.1uF και 10uF όσο το δυνατόν κοντά στους ακροδέκτες. Οι πυκνωτές αυτοί συμβάλλουν στη διασφάλιση σταθερής τροφοδοσίας, καθώς οι πυκνωτές των 0.1u καταστέλλουν κυρίως διακυμάνσεις υψηλής συχνότητας, ενώ οι 10uF λειτουργούν ως τοπικές αποθήκες ενέργειας περιορίζοντας τις διακυμάνσεις χαμηλότερων συχνοτήτων.

Τέλος, ως διεπαφή των ηλεκτροδίων με την πλακέτα επιλέχθηκε ένας απλός διπολικός ακροδέκτης (2-terminal connector) με βήμα 5mm. Στο Σχήμα 4.13 φαίνεται η τελική πλακέτα με όλα τα στοιχεία ενσωματωμένα.


\begin{figure}[H]
  \centering
  \includegraphics[width=0.85\textwidth]{soil.jpg}
  \caption{Η τελική μορφή της πλακέτας.}
  \label{fig:signal-buffering-filtering}
\end{figure}
